\chapter{Modular Invariance and the Weyl Anomaly}
\label{ch:volume-iii-modular-invariance-weyl-anomaly}

\section{Torus Partition Function}

Put a two-dimensional CFT on a Euclidean torus with complex modulus
\[
  \tau\in\mathbb H,
  \qquad
  q=\ee^{2\pi\ii\tau}.
\]
The Hilbert space on the spatial circle has Hamiltonian
\[
  H_{\rm cyl}
  =
  L_0-\frac c{24}
  +
  \widetilde L_0-\frac{\widetilde c}{24},
\]
and momentum
\[
  P_{\rm cyl}
  =
  L_0-\frac c{24}
  -
  \left(\widetilde L_0-\frac{\widetilde c}{24}\right).
\]
The torus partition function is
\begin{equation}
  Z(\tau,\bar\tau)
  =
  \operatorname{Tr}_{\Hilb_{S^1}}
  q^{\,L_0-c/24}
  \bar q^{\,\widetilde L_0-\widetilde c/24}.
  \label{eq:torus-partition-function}
\end{equation}
This expression is the thermal trace in one channel and the Euclidean path
integral on the torus in geometric language.

\section{Modular Transformations}

The torus can be presented as
\[
  \C/(\mathbb Z+\tau\mathbb Z).
\]
Changing the basis of cycles by a matrix
\[
  \begin{pmatrix}
  a & b\\
  c & d
  \end{pmatrix}
  \in SL(2,\mathbb Z)
\]
changes the modulus by
\begin{equation}
  \tau\mapsto\frac{a\tau+b}{c\tau+d}.
  \label{eq:modular-action-tau}
\end{equation}
The generators may be taken as
\[
  T:\tau\mapsto\tau+1,
  \qquad
  S:\tau\mapsto-\frac1\tau .
\]
Locality and invariance of the path integral under diffeomorphisms of the
torus require
\begin{equation}
  Z(\tau,\bar\tau)
  =
  Z\left(\frac{a\tau+b}{c\tau+d},
          \frac{a\bar\tau+b}{c\bar\tau+d}\right).
  \label{eq:modular-invariance-partition-function}
\end{equation}
The \(S\) transformation exchanges the thermal and spatial cycles; it equates
the low-temperature and high-temperature expansions of the same partition
function.

\begin{figure}[H]
  \centering
  \resizebox{0.92\textwidth}{!}{%
  \begin{tikzpicture}[x=1cm,y=1cm,every node/.style={font=\scriptsize}]
    \tikzset{arrow/.style={-{Latex[length=2mm]},line width=0.8pt}}
    \begin{scope}
      \node[font=\small] at (0,2.0) {thermal cylinder};
      \draw[thick] (-0.95,1.20) ellipse (0.75 and 0.18);
      \draw[thick] (-1.70,1.20)--(-1.70,-1.20);
      \draw[thick] (-0.20,1.20)--(-0.20,-1.20);
      \draw[thick] (-0.95,-1.20) ellipse (0.75 and 0.18);
      \draw[arrow,blue!65!black] (-2.05,-1.10)--(-2.05,1.10);
      \node[blue!65!black,left] at (-2.05,0) {thermal cycle};
      \draw[arrow,orange!85!black] (-1.70,-1.55)--(-0.20,-1.55);
      \node[orange!85!black,below] at (-0.95,-1.55) {spatial circle};
    \end{scope}

    \draw[arrow] (1.10,0)--(2.30,0);
    \node[align=center] at (1.70,0.50) {identify\\ends};

    \begin{scope}[shift={(4.15,0)}]
      \node[font=\small] at (0,2.0) {torus lattice};
      \draw[thick] (-1.25,-0.90)--(1.15,-0.90)--(1.65,0.90)--(-0.75,0.90)--cycle;
      \draw[blue!65!black,line width=0.95pt] (-1.25,-0.90)--(1.15,-0.90);
      \draw[orange!85!black,line width=0.95pt] (-1.25,-0.90)--(-0.75,0.90);
      \node[blue!65!black] at (-0.05,-1.23) {\(1\)};
      \node[orange!85!black] at (-1.35,0.05) {\(\tau\)};
    \end{scope}

    \draw[arrow] (6.45,0)--(7.55,0);
    \node[align=center] at (7.00,0.50) {change\\cycle basis};

    \begin{scope}[shift={(9.35,0)}]
      \draw[thick] (-1.15,-0.85)--(1.15,-0.85)--(0.70,0.85)--(-1.60,0.85)--cycle;
      \draw[violet,line width=0.95pt] (-1.15,-0.85)--(0.70,0.85);
      \draw[green!45!black,line width=0.95pt] (-1.60,0.85)--(1.15,-0.85);
      \node[violet] at (0.05,1.15) {\(S\)};
      \node[green!45!black] at (0,-1.20) {\(T\)};
    \end{scope}
  \end{tikzpicture}
  }
  \caption{The same torus admits different choices of spatial and thermal
  cycles.  Modular invariance is invariance under this change of cycle basis.}
  \label{fig:torus-modular-cycles}
\end{figure}

\section{Characters and Spectrum}

If the Hilbert space decomposes into irreducible chiral representations
\(\mathcal V_i\otimes\widetilde{\mathcal V}_j\) with multiplicities
\(N_{ij}\), then
\[
  Z(\tau,\bar\tau)
  =
  \sum_{i,j}N_{ij}\chi_i(\tau)\overline{\chi_j(\tau)}.
\]
Here
\[
  \chi_i(\tau)
  =
  \operatorname{Tr}_{\mathcal V_i}
  q^{L_0-c/24}
\]
is the chiral character.  In rational CFTs, the characters span a
finite-dimensional representation of \(SL(2,\mathbb Z)\):
\[
  \chi_i(-1/\tau)=\sum_j S_{ij}\chi_j(\tau),
  \qquad
  \chi_i(\tau+1)=\sum_j T_{ij}\chi_j(\tau).
\]
The matrix \(N_{ij}\) must commute with these \(S\) and \(T\) matrices.  Thus
modular invariance is a finite algebraic constraint in rational CFT and a
spectral constraint in general.

\section{Weyl Anomaly}

Let
\[
  W[g]=-\log Z[g]
\]
be the Euclidean effective action on a curved two-dimensional background
\((\Sigma,g)\).  Under an infinitesimal Weyl transformation
\[
  \delta_\omega g_{\mu\nu}=2\omega g_{\mu\nu},
\]
the anomaly is
\begin{equation}
  \delta_\omega W[g]
  =
  -\frac{c}{24\pi}
  \int_\Sigma\dd^2x\sqrt g\,\omega R.
  \label{eq:weyl-anomaly-variation}
\end{equation}
With the convention
\[
  \delta W
  =
  \frac12
  \int\dd^2x\sqrt g\,
  \langle T^{\mu\nu}\rangle\,\delta g_{\mu\nu},
\]
this is equivalent to
\begin{equation}
  \langle T^\mu{}_\mu\rangle_g
  =
  -\frac{c}{24\pi}R.
  \label{eq:trace-anomaly-2d}
\end{equation}

For a finite Weyl transformation
\[
  g'_{\mu\nu}=\ee^{2\omega}g_{\mu\nu},
\]
the partition function transforms by the Wess--Zumino action
\begin{equation}
  Z[g']
  =
  \exp\!\left(-S_{\rm WZ}[\omega;g]\right)Z[g],
  \label{eq:finite-weyl-transform-partition}
\end{equation}
where
\begin{equation}
  S_{\rm WZ}[\omega;g]
  =
  \frac{c}{24\pi}
  \int_\Sigma\dd^2x\sqrt g\,
  \left(
    g^{\mu\nu}\partial_\mu\omega\partial_\nu\omega
    +
    R\omega
  \right).
  \label{eq:wess-zumino-weyl-action}
\end{equation}
The cocycle property of this functional is the finite form of the anomaly.

\section{From Local Data to Surfaces}

Every two-dimensional metric is locally conformally flat.  Local correlators
on a curved Riemann surface are therefore obtained from plane correlators,
coordinate transition functions, and the Weyl-anomaly factor.  The remaining
global data are the moduli of the surface and the choice of sewing channel.

\begin{figure}[H]
  \centering
  \resizebox{0.90\textwidth}{!}{%
  \begin{tikzpicture}[x=1cm,y=1cm,every node/.style={font=\scriptsize}]
    \tikzset{
      insertion/.style={circle,fill=violet!80!black,inner sep=1.7pt},
      arrow/.style={-{Latex[length=2mm]},line width=0.8pt}
    }
    \draw[thick]
      (-4.65,0.25) .. controls (-4.20,1.05) and (-3.20,1.05) .. (-2.70,0.32)
      .. controls (-2.20,-0.45) and (-1.30,-0.35) .. (-0.85,0.25)
      .. controls (-0.45,-0.82) and (-1.65,-1.10) .. (-2.55,-0.75)
      .. controls (-3.50,-0.42) and (-4.25,-0.85) .. (-4.65,0.25);
    \node[insertion] at (-4.05,0.22) {};
    \node[insertion] at (-1.50,-0.18) {};
    \node[violet!80!black] at (-4.25,0.60) {\(p_i\)};
    \node[violet!80!black] at (-1.25,0.18) {\(p_j\)};
    \node at (-2.75,-1.55) {curved surface};

    \draw[arrow] (-0.10,0)--(1.05,0);
    \node[align=center] at (0.48,0.52) {local\\Weyl chart};

    \draw[thick] (1.75,-1.00) rectangle (5.20,1.00);
    \foreach \x/\y/\lab in {2.35/0.40/{z_i},4.45/-0.30/{z_j}} {
      \node[insertion] at (\x,\y) {};
      \node[violet!80!black,above right] at (\x,\y) {\(\lab\)};
    }
    \draw[green!45!black,line width=0.85pt] (2.35,0.40) circle (0.35);
    \draw[green!45!black,line width=0.85pt] (4.45,-0.30) circle (0.35);
    \node[blue!65!black] at (3.45,-1.35)
    {plane OPE data plus anomaly};
  \end{tikzpicture}
  }
  \caption{On a Riemann surface, local flat coordinates reduce local
  correlators to plane data, while Weyl transformations contribute the anomaly
  functional.}
  \label{fig:weyl-local-chart}
\end{figure}

Cutting a Riemann surface into three-holed spheres expresses its correlator in
terms of three-point data and sums over intermediate states.  Different
decompositions of the same surface are related by crossing moves and modular
moves.  Thus four-point crossing and torus modular invariance are local
manifestations of the same global consistency requirement.
